% A structured grid in 3D: the cells are identical and their neighbours follow
% from the indices, so no connectivity has to be stored. Drawn isometrically by
% projecting the three axes onto the page.
\documentclass[dvisvgm,border=3pt]{standalone}
\input{preamble.tex}
\begin{document}
\begin{tikzpicture}[x=1cm,y=1cm]

  \def\N{4}
  \def\s{0.52}                 % cell size
  % isometric axes: x to the right, y into the page, z up
  \newcommand{\px}[3]{{(#1)*\s + (#2)*\s*0.52}}
  \newcommand{\py}[3]{{(#3)*\s + (#2)*\s*0.30}}

  % \face{corner coords as four (i,j,k) triples}{fill}
  \newcommand{\quadf}[9]{%
    \fill[#9,draw=inkdark,line width=0.4pt]
      (\px{#1}{#2}{#3},\py{#1}{#2}{#3}) -- (\px{#4}{#5}{#6},\py{#4}{#5}{#6})
      -- (\px{#7}{#8}{0},\py{#7}{#8}{0}) -- cycle;}

  % --- top face (z = N) ---------------------------------------------------
  \foreach \i in {0,...,3}{\foreach \j in {0,...,3}{
    \pgfmathsetmacro{\t}{Mod(\i+\j,2) ? 100 : 45}
    \fill[uibkorange!\t,draw=inkdark,line width=0.4pt]
      (\px{\i}{\j}{\N},\py{\i}{\j}{\N})
      -- (\px{\i+1}{\j}{\N},\py{\i+1}{\j}{\N})
      -- (\px{\i+1}{\j+1}{\N},\py{\i+1}{\j+1}{\N})
      -- (\px{\i}{\j+1}{\N},\py{\i}{\j+1}{\N}) -- cycle;}}

  % --- front face (y = 0) -------------------------------------------------
  \foreach \i in {0,...,3}{\foreach \k in {0,...,3}{
    \pgfmathsetmacro{\t}{Mod(\i+\k,2) ? 100 : 45}
    \fill[uibkorange!\t,draw=inkdark,line width=0.4pt]
      (\px{\i}{0}{\k},\py{\i}{0}{\k})
      -- (\px{\i+1}{0}{\k},\py{\i+1}{0}{\k})
      -- (\px{\i+1}{0}{\k+1},\py{\i+1}{0}{\k+1})
      -- (\px{\i}{0}{\k+1},\py{\i}{0}{\k+1}) -- cycle;}}

  % --- right face (x = N) -------------------------------------------------
  \foreach \j in {0,...,3}{\foreach \k in {0,...,3}{
    \pgfmathsetmacro{\t}{Mod(\j+\k,2) ? 70 : 30}
    \fill[uibkorange!\t,draw=inkdark,line width=0.4pt]
      (\px{\N}{\j}{\k},\py{\N}{\j}{\k})
      -- (\px{\N}{\j+1}{\k},\py{\N}{\j+1}{\k})
      -- (\px{\N}{\j+1}{\k+1},\py{\N}{\j+1}{\k+1})
      -- (\px{\N}{\j}{\k+1},\py{\N}{\j}{\k+1}) -- cycle;}}

\end{tikzpicture}
\end{document}
